\section{Въведение и мотивация за курса}

Този курс има за цел да изгради у вас така нареченото стохастично (вероятностно) мислене. Въпреки че курсът е чисто математически и ще изисква усвояване на строги дефиниции и теоретичен апарат, този поглед към света е изключително полезен и намира приложение на най-високо ниво във фундаменталната наука, в съвременните технологии (като изкуствен интелект и машинно обучение), както и в напълно ежедневни, злободневни ситуации.

\subsection{Примери за приложение на стохастиката}

За да се мотивира необходимостта от такъв апарат, ще разгледаме няколко примера, които показват мощта на вероятностните модели.

\textbf{Брауново движение.} През 1827 г. ботаникът Робърт Браун (Р. Браун 1827; 1905) наблюдава хаотичното движение на частица полен, поставена на водна повърхност. На онзи етап от развитието на науката атомарната теория все още не е била утвърдена и не е било ясно дали това движение се дължи на вътрешна енергия на частицата, или на външната среда. Едва през 1905 г. Алберт Айнщайн успява да опише математически това движение, стъпвайки на кинетичната теория. Резултатът постулира, че в резултат от ударите на водните молекули от всички посоки, частицата приема хаотична траектория, която постоянно мени посоката си поради моментния дисбаланс на силите.
Любопитно е, че ако разгледаме Брауновото движение като чисто стохастичен обект, траекторията е непрекъсната, но никъде диференцируема и има безкрайна дължина. Въпреки тази физическа „нерационалност“, моделът е изключително добро приближение на реалността. Дори днес стохастичните диференциални уравнения са в основата на моделирането на сложни физични и финансови феномени.


\textbf{Тотализаторът и шансът.} На 6 септември 2009 г. в България се случва събитие, което предизвиква огромен медиен скандал: в два последователни тиража на тотализатора „6 от 49“ се падат едни и същи числа. Примерът е взет от книга на колеги от Imperial College, които го използват, за да покажат как дори елементарни вероятностни разсъждения помагат при анализа на реални събития. Обикновено първосигналната реакция е да се търси измама и манипулация. Въпреки това, ако се направи елементарен статистически анализ (взимайки предвид, че до онзи момент са изиграни около 10 000 тиража), се оказва, че вероятността за подобно събитие, макар и малка, съвсем не е пренебрежима. Разбирането на тези базови вероятности ни позволява да оценяваме обективно фактите, вместо да се поддаваме на емоционални спекулации. Допълнителен аргумент против хипотезата за измама е фактът, че във втория тираж печалбата е била разделена между 14 души — едва ли някой би организирал измама, за да дели печалбата с още толкова хора.


\textbf{Статистика и актуални проблеми.} В контекста на пандемията от COVID-19 се появяват множество спекулации. Например, когато се съобщи, че около 4-5\% от починалите за деня са били ваксинирани, веднага започват конспиративни теории за фалшиви сертификати или за пълна неефективност на ваксините. Отново, един елементарен статистически анализ на базата на съотношението между ваксинирани и неваксинирани в обществото би показал реалната ефективност на ваксините, която, макар и не 100\%, съществува.


\textbf{Фундаментална математика.} Преди няколко години известният математик Терънс Тао доказа в голяма степен хипотезата на българския академик Благовест Сендов. Ключов момент от неговата изключително сложна работа отново беше базиран на стохастиката.


\section{Елементарни събития и пространство от елементарни събития}

За да формализираме тези идеи и да започнем да градим математическата теория, трябва да въведем основните понятия, свързани със случайните явления.

\begin{defn}
Под \emph{случаен експеримент} разбираме изследване, опит или наблюдение, чийто точен изход не е предварително известен (не е предопределен). Например, при хвърляне на монета не знаем дали ще се падне „ези“ или „тура“, защото не разполагаме с цялата необходима информация (сили, въздушно съпротивление, начален импулс и др.).
\end{defn}

\begin{defn}
Възможните изходи от даден случаен експеримент се наричат \emph{елементарни събития} и обикновено се обозначават с малката гръцка буква $\omega$.
\end{defn}

\begin{defn}
За даден случаен експеримент, множеството от всички възможни елементарни събития се нарича \emph{пространство от елементарни събития} и се означава с главно $\Omega$.
\end{defn}

Нека разгледаме няколко конкретни примера за пространства от елементарни събития:

\textbf{Хвърляне на монета.}
Експеримент с бинарен изход.
\[
\Omega = \{ \text{'ези'}; \text{'тура'} \}
\]
За удобство и компютърна обработка, много често тези два изхода се кодират с 0 и 1:
\[
\Omega = \{ 0; 1 \}
\]


\textbf{Хвърляне на зарче.}
Експериментът се състои в хвърляне на стандартен шестстенен зар, независимо дали е честен или не. Възможните изходи са броят точки на горната стена:
\[
\Omega = \{ 1, 2, 3, 4, 5, 6 \}
\]
Едно конкретно елементарно събитие, например падането на числото 5, се записва като $\omega_5 = \{ 5 \}$.


\textbf{Един тираж на тото „6 от 49“.}
Тук експериментът е изтеглянето на шест числа от общо 49. Всяко елементарно събитие е една конкретна (неподредена) шесторка от числа.
\[
\Omega = \{ \omega_1, \omega_2, \dots, \omega_N \}
\]
Броят на възможните изходи (кардиналността на множеството) е комбинация от 49 елемента 6-ти клас:
\[
N = \binom{49}{6} = \underline{\underline{13983816}}
\]


\textbf{Неизброимо множество — време на живот.}
Да си представим, че изследваме времето на живот на един процесор. То може да бъде всяко неотрицателно реално число (дори 0, ако е фабрично дефектен). В този случай пространството е непрекъснато (неизброимо):
\[
\Omega = \mathbb{R}_+ = [0, \infty)
\]


\textbf{Брауново движение.}
За полена върху водната повърхност, $\Omega$ се състои от всички възможни непрекъснати криви (функции), които стартират от началото на координатната система и описват траектория. Това е изключително сложно и голямо функционално пространство. За щастие, в повечето практически задачи не е нужно да описваме експлицитно цялата тази структура.


\section{Събития и операции с тях}

В реалните проблеми рядко се интересуваме от това дали се е случило точно едно конкретно елементарно събитие. Много по-често въпросите ни касаят \emph{група} от елементарни събития, които споделят общ признак.

\begin{defn}
Всяко подмножество $A$ на пространството от елементарни събития $\Omega$ ($A \subseteq \Omega$) се нарича \emph{събитие}.
\end{defn}

Например, ако разглеждаме клиентите, влизащи в един магазин, $\Omega$ може да съдържа всички възможни хора. Събитието $A = \{\text{жени}\}$ обединява всички елементарни събития, които отговарят на признака пол. Магазинът обикновено не се интересува дали точно Иван или Деница е влязъл, а по-скоро от общите характеристики на клиентите (напр. мъж/жена).
Ако хвърляме зарче, събитието $A = \{ 1, 3, 5 \}$ представлява падането на нечетно число. Групирали сме изходите по признака „нечетност“.

\subsection{Операции със събития}

Тъй като събитията са множества, с тях можем да извършваме стандартните теоретико-множествени операции, които имат ясен вероятностен смисъл. Нека $A, B \subseteq \Omega$.

\textbf{Влагане (Подмножество):}
\[
A \subseteq B
\]
Казваме, че $A$ е подмножество на $B$, ако настъпването на $A$ гарантира настъпването на $B$. Формално: ако $\omega \in A \implies \omega \in B$. (Събитията са равни, $A = B$, ако $A \subseteq B$ и $B \subseteq A$).

\textbf{Обединение (Логическо „ИЛИ“):}
\[
C = A \cup B
\]
Това е събитието $C$, такова че ако $\omega \in C$, то $\omega \in A$ \textbf{или} $\omega \in B$.
Например, ако $A = \{\text{гласували за партия } X\}$, а $B = \{\text{хора на възраст 19-21 год.}\}$, то обединението им $A \cup B$ съдържа всички избиратели, които са или гласували за партия $X$, или са на възраст между 19 и 21 години (или и двете едновременно).

\textbf{Сечение (Логическо „И“):}
\[
C = A \cap B
\]
Това е събитието $C$, такова че ако $\omega \in C$, то $\omega \in A$ \textbf{и} $\omega \in B$.
В същия пример, сечението $A \cap B$ ще бъдат точно онези избиратели на възраст между 19 и 21 години, които са гласували за партия $X$. За една партия е важно да анализира тези сечения, за да знае какви са исканията на тази конкретна демографска група.

\textbf{Допълнение (Отрицание):}
\[
A^c = \Omega \setminus A
\]
Събитието $A^c$ се случва тогава и само тогава, когато събитието $A$ не се случва. Формално: ако $\omega \in A^c$, то $\omega \notin A$.
Разбира се, в сила е:
\[
\Omega = A \cup A^c
\]
Примери за допълнение:
Ако $A = \{\text{жени}\}$, то $A^c = \{\text{мъже}\}$.
Ако $A = \{1, 3, 5\}$, то $A^c = \{2, 4, 6\}$.
Ако $A = \{\text{гласували за } X\}$, то $A^c = \{\text{негласували за } X\}$.

\subsection{Свойства на операциите}

Свойствата на тези операции са интуитивни и добре познати от дискретната математика:

\textbf{а) Комутативност:}
\[
A \cup B = B \cup A; \quad A \cap B = B \cap A
\]

\textbf{б) Асоциативност:}
\[
A \cup (B \cup C) = (A \cup B) \cup C = A \cup B \cup C
\]
\[
A \cap (B \cap C) = (A \cap B) \cap C = A \cap B \cap C
\]

\textbf{в) Дистрибутивност:}
\[
A \cup (B \cap C) = (A \cup B) \cap (A \cup C)
\]

\textbf{г) Инволюция на допълнението:}
Ако вземем допълнението на допълнението, се връщаме в изходното множество:
\[
(A^c)^c = A
\]

\textbf{д) Безкрайни обединения и сечения:}
Ако $A_i$ за $i \ge 1$ е редица от събития (изброимо много), можем да дефинираме тяхното обединение и сечение:
\[
\bigcup_{i=1}^{\infty} A_i = \{ \omega \in \Omega : \omega \text{ принадлежи на поне едно } A_i \}
\]
\[
\bigcap_{i=1}^{\infty} A_i = \{ \omega \in \Omega : \omega \text{ принадлежи на } \forall A_i \}
\]
Тези безкрайни операции са фундаментални за аксиоматичното изграждане на теорията на вероятностите (подобно на границите в математическия анализ), въпреки че в практиката най-често работим с крайни множества.

\textbf{е) Закони на Де Морган (De Morgan):}
Законите на Де Морган свързват обединението, сечението и допълнението:
За две събития:
\[
(A \cup B)^c = A^c \cap B^c
\]
За изброимо много събития:
\[
\left(\bigcup_{i=1}^\infty A_i\right)^c = \bigcap_{i=1}^\infty A_i^c
\]
\[
\left(\bigcap_{i=1}^{\infty} A_i\right)^c = \bigcup_{i=1}^{\infty} A_i^c
\]
С думи: допълнението на обединението е сечение от допълненията, а допълнението на сечението е обединение от допълненията.

\section{Сигма-алгебра ($\sigma$-алгебра)}

За да дефинираме коректно вероятностна мярка, не винаги можем просто да разглеждаме всички възможни подмножества на $\Omega$. Необходима ни е математическа структура, която да съдържа събитията, които можем да „измерим“, и тази структура трябва да бъде затворена относно определени операции (за да не се окаже, че ако знаем вероятността на събитията $A$ и $B$, не можем да сметнем вероятността на тяхното обединение). Тази структура се нарича $\sigma$-алгебра.

\begin{keydefn}[$\sigma$-алгебра]\label{def:sigma-algebra}
Нека $\Omega$ е пространство от елементарни събития. Нека $\mathcal{A}$ е колекция от подмножества на $\Omega$. Тогава $\mathcal{A}$ се нарича \emph{$\sigma$-алгебра}, ако са изпълнени следните три условия:
\begin{itemize}
    \item[(i)] $\emptyset \in \mathcal{A}$ \quad (Празното множество е част от колекцията.)
    \item[(ii)] Ако $A \in \mathcal{A}$, то $A^c \in \mathcal{A}$ \quad (Колекцията е затворена относно операцията допълнение.)
    \item[(iii)] Ако $A_i \in \mathcal{A}$ за $\forall i \ge 1$, то $\bigcup_{i=1}^{\infty} A_i \in \mathcal{A}$ \quad (Колекцията е затворена относно изброими обединения.)
\end{itemize}
\end{keydefn}

Затвореността за обединение на две множества $A \cup B \in \mathcal{A}$ следва от (iii), като приемем, че останалите множества от безкрайната редица са празни.

От тези три аксиоми можем да изведем важно следствие за операцията сечение, която не фигурира експлицитно в дефиницията.

\begin{cor}
Нека $\mathcal{A}$ е $\sigma$-алгебра и $A_i \in \mathcal{A}$ за $\forall i \ge 1$. Тогава:
\[
B = \bigcap_{i=1}^{\infty} A_i \in \mathcal{A}
\]
Тоест, $\sigma$-алгебрата е затворена и относно изброими сечения.
\end{cor}

\begin{proof}
Искаме да покажем, че $\mathcal{A} \ni B = \bigcap_{i=1}^{\infty} A_i$.
От свойство (ii) знаем, че ако допълнението $B^c \in \mathcal{A}$, то и самото множество $B = (B^c)^c \in \mathcal{A}$. Затова ще изследваме допълнението на $B$. Използвайки законите на Де Морган, обръщаме сечението в обединение:
\[
B^c = \left( \bigcap_{i=1}^{\infty} A_i \right)^c = \bigcup_{i=1}^{\infty} A_i^c
\]
От свойство (ii) на $\sigma$-алгебрата, тъй като $A_i \in \mathcal{A}$, следва че и $A_i^c \in \mathcal{A}$.
Тогава, съгласно свойство (iii), изброимото обединение на множества, които принадлежат на $\mathcal{A}$, също принадлежи на $\mathcal{A}$. Следователно:
\[
\bigcup_{i=1}^{\infty} \underbrace{A_i^c}_{\substack{\in \mathcal{A} \\ \text{от (ii)}}} \stackrel{\text{(iii)}}{\in} \mathcal{A}
\]
Така показахме, че $B^c \in \mathcal{A}$. Прилагайки отново свойство (ii), заключаваме, че $B = (B^c)^c \in \mathcal{A}$. С това следствието е доказано.
\end{proof}

\subsection{Примери за $\sigma$-алгебри}

\begin{example}[Тривиална $\sigma$-алгебра]\label{ex:01-10}
За експеримент с $\Omega = \{0, 1\}$ (хвърляне на монета), най-простата $\sigma$-алгебра е:
\[
\mathcal{A} = \{\emptyset, \Omega\}
\]
Тя отговаря на условията, но е безполезна за практически изчисления, тъй като не съдържа информация за самите изходи (0 или 1).
\end{example}

\begin{example}[Множество от всички подмножества: Булеан]\label{ex:01-11}
За същото $\Omega = \{0, 1\}$, естествената $\sigma$-алгебра е колекцията от всички възможни подмножества:
\[
\mathcal{A} = 2^\Omega = \{\emptyset, \Omega, \{0\}, \{1\}\}
\]
Ако $\Omega$ има $n$ елемента, то $\mathcal{A} = 2^\Omega$ има $2^n$ елемента. За дискретни, крайни или изброимо безкрайни пространства, обикновено работим точно с тази най-богата $\sigma$-алгебра, защото върху нея винаги може коректно да се дефинира вероятност.
\end{example}

\begin{example}[Непрекъснати пространства и Борелова $\sigma$-алгебра]\label{ex:01-12}
Когато разглеждаме неизброими пространства, например $\Omega = \mathbb{R}_+ = [0, \infty)$, колекцията от всички подмножества $2^\Omega$ става „твърде голяма“. Математически е невъзможно да се зададе вероятностна мярка, която да измерва \emph{всяко едно} подмножество на реалните числа. За щастие, в практиката ние никога не се интересуваме от екзотични и „изкривени“ множества. Вълнуваме се от интервали от вида $A = (a, b)$ — например дали температурата е между 10 и 12 градуса. $\sigma$-алгебрата, породена от всички такива интервали, се нарича \emph{Борелова $\sigma$-алгебра} и върху нея може коректно да се зададе вероятност.
\end{example}

\begin{example}[Редуциране на неизброимо пространство]\label{ex:01-13}
Представете си, че стреляме с лък по мишена за дартс. Точната точка на попадение е от неизброимото множество на двумерната равнина. Ние обаче се интересуваме само от това колко точки ще получим, т.е. в кой цветен сектор е попаднала стрелата (например: 1 — черно, 2 — синьо, 3 — червено). Така ние свеждаме сложния неизброим експеримент до просто пространство с три елементарни изхода:
\[
\Omega = \{\omega_1, \omega_2, \omega_3\}
\]
Съответната $\sigma$-алгебра ще бъде $2^\Omega$ и ще съдържа точно $2^3 = 8$ елемента:
\[
\mathcal{A} = 2^\Omega = \{\emptyset, \Omega, \{\omega_1\}, \{\omega_2\}, \{\omega_3\}, \{\omega_1, \omega_2\}, \{\omega_1, \omega_3\}, \{\omega_2, \omega_3\}\}
\]
Това улеснява значително анализа и елиминира нуждата от сложен математически апарат там, където практическият проблем не го изисква.
\end{example}

\section{Математическо моделиране на проблема с тотализатора}

Нека се върнем на медийния скандал с тотализатора „6 от 49“. Искаме да моделираме събитието: \emph{„Паднали са се едни и същи числа в два последователни тиража в рамките на 10 000 изиграни тиража.“}

Пространството от елементарни събития за всички 10 000 тиража заедно може да се опише като Декартово произведение на индивидуалните пространства на всеки тираж:
\[
\Omega = \bigotimes_{i=1}^{10000} \Omega_i
\]
Едно конкретно елементарно събитие (история на всички тегления) изглежда така:
\[
\omega = (\omega^{(1)}, \omega^{(2)}, \dots, \omega^{(10000)})
\]
където всяко $\omega^{(i)}$ е шесторка числа от $i$-тия тираж.
Мощността на това пространство е колосална:
\[
|\Omega| = (\underbrace{13983816}_{N})^{10000} = N^{10000}
\]
Съответната $\sigma$-алгебра $\mathcal{A} = 2^\Omega$ ще има $2^{N^{10000}}$ елемента, което е необозримо голямо число, но важното е, че е \emph{крайно}, следователно нямаме теоретични проблеми с измерването.

За да пресметнем вероятността на разглежданото скандално събитие $A$, ние го разбиваме на по-прости, елементарни компоненти. Дефинираме $A_i$ като събитието: \emph{„шесторката в $i$-тия тираж е същата като шесторката в $(i+1)$-вия тираж“}.
Формално:
\[
A_i = \{ \omega^{(i)} = \omega^{(i+1)} \} \subseteq \Omega, \quad A_i \in 2^\Omega
\]
Тогава търсеното събитие $A$ („едни и същи числа в два последователни тиража от всички 10 000“) е просто обединението на всички възможни такива съвпадения за $i$ от 1 до 9999:
\[
A = \bigcup_{i=1}^{9999} A_i
\]
Превеждайки словесното описание на събитието в стриктна теоретико-множествена форма, ние създаваме солидна основа за пресмятане на неговата вероятност. В следващите лекции ще видим как да оценим вероятността на подобно обединение. Макар тя да не е голяма, оценката показва, че не е толкова малка, че сама по себе си да предизвика наистина сериозно съмнение.

\section{Задачи}

\begin{enumerate}
\item Проверете закона на Де Морган за две множества,
\[
(A \cup B)^c = A^c \cap B^c,
\]
а по желание проверете и другия закон на Де Морган.

\item \textbf{Парадоксът с двата плика.}

В края на лекцията оставяме една класическа, интригуваща задача.
Представете си, че имате два затворени плика. Във всеки плик има лист с написано реално число, съответно $a$ и $b$, като $a < b$ (двете числа са различни). Организаторът на играта знае числата, но вие нямате абсолютно никаква предварителна (априорна) информация в какъв диапазон са те (може да са 1 и 2, може да са 1 и 10 000, може да са отрицателни).
Избирате случайно един от пликовете и го отваряте. Виждате числото в него. Целта ви е да познаете дали това е по-голямото число, или по-малкото (или респективно да решите дали да смените плика, за да вземете по-голямата сума).

\textbf{Въпросът е:} Съществува ли стратегия, която да ви гарантира вероятност строго по-голяма от $\frac{1}{2}$ (тоест по-добра от чисто налучкване), че ще познаете/изберете по-голямото число?
Помислете върху този проблем, без да спекулирате произволно, а търсейки строга математическа логика.

\begin{supp}[Решение на задачата с двата плика]\label{supp:envelopes}
Отговорът е \emph{да}, и ключът е стратегията да бъде \emph{случайна}.

Нека $C$ е събитието „сменям плика“, а $A$ и $B$ — съответно събитията „видял съм
по-малкото число $a$“ и „видял съм по-голямото число $b$“, всяко с вероятност $\tfrac12$.
Печелим, ако при $A$ сменим, а при $B$ не сменим, тоест вероятността за успех е
\[
\mathbb{P}(\text{успех})
 = \tfrac12\,\mathbb{P}(C \given A) + \tfrac12\big(1 - \mathbb{P}(C \given B)\big)
 = \tfrac12 + \tfrac12\big(\mathbb{P}(C \given A) - \mathbb{P}(C \given B)\big).
\]
Следователно всяка стратегия, при която сменяме \emph{по-охотно} след малко число,
отколкото след голямо — тоест $\mathbb{P}(C \given A) > \mathbb{P}(C \given B)$ — печели с
вероятност строго над $\tfrac12$.

Такава стратегия съществува, при това без никакво предварително знание за $a$ и $b$:
след като видим число $x$, сменяме с вероятност $e^{-x}$ (или с която и да е строго
намаляваща функция със стойности в $(0,1)$). Тогава
\[
\mathbb{P}(C \given A) - \mathbb{P}(C \given B) = e^{-a} - e^{-b} > 0,
\]
понеже $a < b$. Печалбата е малка, ако $a$ и $b$ са близки, но е строго положителна винаги.
\end{supp}
\end{enumerate}
